% !PRINTED none

\begin{center}
\gesperrt{Achter Abschnitt.}

\textbf{\gesperrt{Die Poly\"edergruppen.}}

\rule{3cm}{0.4pt}

\S.\ 55.
\end{center}

\begin{center}
\gesperrt{Die cyklischen Gruppen und die Di\"edergruppen.}
\end{center}

Wir gehen nun dazu über, die allgemeinen Principien zur
wirklichen Bildung der verschiedenen Poly\"edergruppen anzuwenden,
zunächst also zu zeigen, dass die in \S.\,52 als möglich
erkannten Arten dieser Gruppen alle existiren.

Bei den cyklischen Gruppen $n^{\text{ten}}$ Grades haben wir nur zwei
Systeme conjugirter Pole, deren jedes nur einen $(n-1)$fachen
Pol enthält. Diese beiden Pole müssen also die gemeinsamen
Pole aller Substitutionen der Gruppe sein und können nach
\S.\,51, 2. als $0$ und $\infty$ angenommen werden. Wir haben dann
nur die beiden linearen Grundformen
\weq{1}{f_1 = x_1, \quad f_2 = x_2.}
Alle Substitutionen der Gruppe sind multiplicativ, und der
Multiplicator muss eine $n^{\text{te}}$ Einheitswurzel sein. Sie müssen also
die Form haben:
\weq{2}{x, \varepsilon x, \varepsilon^2 x, \ldots, \varepsilon^{n-1} x,}
wenn $\varepsilon$ eine primitive $n^{\text{te}}$ Einheitswurzel ist.

Wir erhalten also in der That für jedes $n$ eine cyklische
Gruppe, die wir mit $C_n$ bezeichnen wollen:
\weq{3}{C_n = \begin{pmatrix} \varepsilon^r, & 0 \\ 0, & 1 \end{pmatrix}, \quad r = 0, 1, \ldots, n-1,}
oder mit der Determinante $1$ geschrieben:
\weq{4}{C_n = \begin{pmatrix} \varepsilon^{1/2\,r}, & 0 \\ 0, & \varepsilon^{-1/2\,r} \end{pmatrix}.}
